\sectioncentered*{Введение}
\addcontentsline{toc}{section}{Введение}

На сегодняшний день одной из главных тенденций в развитии науки, технологий и промышленности является интеллектуализация. Она направлена на создание систем, обслуживающих нужды человека, и постепенно проникает во все сферы жизни: от производства до искусства и развлечений. Вследствие популяризации различных виртуальных ассистентов, чат-ботов, а также развития «умной» техники, интеллектуальные технологии находят широкое применение в быту и уже стали неотъемлемой частью повседневной жизни.

Под интеллектуальной технологией в данном контексте понимается инструментарий, позволяющий решать задачи, традиционно требующие человеческого мышления. Решение таких задач зачастую не поддается строгому алгоритмическому описанию, поскольку реальные процессы характеризуются высокой степенью неопределенности, сложностью взаимосвязей и наличием большого комплекса внешних и внутренних факторов. Нередко проблемы требуют анализа условий «здесь и сейчас», либо же количество параметров, от которых зависит результат, оказывается критически велико. Этим факторам сопутствуют проблемы качества и полноты входных данных, их динамичность и многомерность. К характерным чертам интеллектуальных задач относятся необходимость многовариантных решений, способных динамически изменяться, потребность в вероятностном подходе, а также адаптация и обучаемость системы.

На текущем этапе научно-технологического развития передовые разработки направлены именно на этот класс задач. Создаются системы, способные комплексно подходить к достижению поставленной цели, компонентами которых выступают технологии искусственного интеллекта: машинное обучение, глубокое обучение и обработка естественного языка. Последнее направление приобретает особую актуальность в проектировании систем персонального ассистирования, так как позволяет пользователю взаимодействовать с программным обеспечением на естественном языке, независимо от уровня технической подготовки. Это способствует повсеместной интеграции интеллектуальных технологий, автоматизации рутинных процессов и повышению эффективности в различных сферах жизни.

В современном мире широкое распространение получили голосовые и текстовые помощники, которые активно участвуют в повседневной жизни, упрощая бытовые операции и поиск информации. Однако особый интерес представляет применение подобных технологий в сфере образования. Поскольку объем знаний в различных областях стремительно накапливается, образовательный процесс приобретает комплексный характер и требует качественной организации. Это особенно актуально для студентов высших учебных заведений и специалистов, занимающихся самообразованием. Для эффективного освоения материала необходимо не просто потреблять информацию, но и грамотно ее структурировать, декомпозировать (разделять на подзадачи) и планировать время для практической отработки. В этом контексте организация времени, или тайм-менеджмент, становится критически важным навыком. Интеллектуальные технологии способны взять на себя роль организатора, автоматизируя процесс планирования и снижая когнитивную нагрузку на обучающегося.

Вместе с тем, анализ текущего состояния предметной области позволяет выявить ряд существенных противоречий. Во-первых, несмотря на активное внедрение цифровых ассистентов и наличие мощных решений на рынке программного обеспечения для планирования, большинство из них ориентировано на корпоративный сегмент. Существует противоречие между избыточностью инструментов горизонтального планирования, предлагаемых универсальными экосистемами (такими как Google или Microsoft), и отсутствием специализированных, вертикально-ориентированных инструментов декомпозиции учебного процесса, необходимых студентам. Универсальный инструментарий зачастую не отвечает специфическим академическим задачам.

Во-вторых, указанные факторы влияют на доступность программного обеспечения. Наблюдается противоречие между финансовыми возможностями обучающихся и обилием дорогостоящих платных подписок на профессиональные сервисы организации времени. Кроме того, немаловажным фактором является критерий региональной доступности программного обеспечения, так как иностранные решения не всегда могут быть использованы в полной мере в конкретных регионах.

В-третьих, специфика образования подразумевает определенные методики, требующие организации личного времени, однако существующие цифровые решения в вузах (системы управления обучением) обычно направлены на формирование отчетности и удовлетворение нужд учебного заведения, а не конкретного студента. Возникает противоречие между наличием обширного цифрового следа обучения (расписания, методические материалы, репозитории кода) и отсутствием инструментов автоматизированной трансформации этого следа в персональный план действий, что оставляет спрос на персонализированный подход неудовлетворенным.

Данные противоречия порождают проблему исследования, которая заключается в необходимости преодоления трудностей, связанных с интенсивным темпом жизни и информационной перегрузкой обучающихся. Перестройка между строгой школьной системой и гибкими долгосрочными проектами в высшем учебном заведении, поддержание баланса между учебой и личной жизнью, а также борьба с прокрастинацией требуют новых технологических подходов к организации времени.

Объектом исследования являются системы персональной организации времени. Предметом исследования выступают технологии и подходы программной реализации персональных ассистентов, а также методы повышения продуктивности и достижения образовательных целей с их помощью.

Целью работы является оптимизация организации времени учащихся вузов за счет делегирования планирования учебных и личных целей и задач в долгосрочной перспективе интеллектуальному ассистенту. Ожидается, что реализация данной цели позволит снизить когнитивную нагрузку на пользователя, повысить продуктивность и вовлеченность в самообразование, а также минимизировать влияние факторов, провоцирующих стресс и прокрастинацию.

Для достижения поставленной цели необходимо решить следующие задачи:
\begin{itemize}
    \item провести анализ предметной области, существующих аналогов и методов реализации ассистентов организации времени;
    \item cпроектировать архитектуру и функциональные модели персонального ассистента, учитывающего специфику образовательного процесса;
    \item разработать программный продукт — персональный ассистент организации времени обучающихся, реализующий функции интеллектуального планирования и декомпозиции задач.
\end{itemize}


В работе применяются теоретические методы исследования, включающие анализ научно-технической литературы и существующих программных решений, а также эмпирические методы, такие как моделирование информационных процессов, проектирование системной архитектуры и методы программной инженерии.

Научная новизна работы заключается в разработке подхода к автоматизированной декомпозиции долгосрочных учебных целей в конкретные календарные слоты с учетом индивидуальной загрузки и академического контекста пользователя, в отличие от традиционных методов ручного планирования.

Теоретическая значимость исследования состоит в обосновании применения агентного подхода и больших языковых моделей для решения задач персонального образовательного менеджмента. Практическая значимость заключается в создании доступного программного инструмента, позволяющего студентам эффективно организовывать учебный процесс, что способствует повышению академической успеваемости и формированию навыков самоорганизации.


